%%%%%%%%%%%%%%%%%%%%% part3.tex %%%%%%%%%%%%%%%%%%%%%%%%%%%%%%%%%%
% 
% Part III: Secure Distributed Systems
%
%%%%%%%%%%%%%%%%%%%%%%%% Springer Nature %%%%%%%%%%%%%%%%%%%%%%

\begin{partbacktext}
\part{Secure Distributed Systems}
\noindent As AI moves from centralized servers to distributed ecosystems—from mobile devices to across-border clusters—the surface area for attack grows exponentially. This third part of the book addresses the security of data and computation in transit and at rest.

In \textbf{Chapter 7}, we introduce the core architecture of \textbf{Federated Learning} (FL), moving beyond the hype to understand its real-world implementation. \textbf{Chapter 8} explores the cryptographic foundations of \textbf{Secure Aggregation}, using Multi-Party Computation and Homomorphic Encryption to protect gradient updates from a curious server. \textbf{Chapter 9} introduces \textbf{Hardware-Rooted Trust}, specifically the use of Trusted Execution Environments (TEEs) like Intel SGX to protect computation in use. Finally, in \textbf{Chapter 10}, we explore \textbf{Robust Aggregation}, ensuring that the global model survives even when participants are malicious or ``Byzantine.''

\begin{casestudy} 
In the fall of 2025, the LedgerNet DeFi protocol, which processed \$1.2B in weekly volume, claimed total privacy through the use of Hardware-Rooted Trust (TEEs). However, a sophisticated side-channel attack targeting the timing of memory accesses within the enclave allowed adversaries to reconstruct account balances for over 10,000 users. This breach exposed the vulnerability of relying on single layers of hardware, driving the adoption of the multi-layered Secure Distributed Systems analyzed in Part III. 
\end{casestudy} 
\end{partbacktext}
